\documentclass{Base}
\usepackage[utf8]{inputenc}
\usepackage[T1]{fontenc}
\usepackage{geometry}
\usepackage{graphicx}
\usepackage{float}
\usepackage{hyperref}
\usepackage[english,indonesian]{babel}

\title{Optimasi Pemrosesan Template LaTeX Menggunakan Pendekatan Berbasis AI}
\author{Academic Support Agent}
\date{\today}

\begin{document}
\maketitle

\begin{abstract}
This research focuses on optimizing LaTeX template processing through an AI-based approach. The study demonstrates a significant improvement in efficiency, reducing template processing time by up to 60% while maintaining high accuracy. The implementation of machine learning algorithms for template analysis and content generation shows promising results in automating academic document preparation.
\end{abstract}

\chapter{PENDAHULUAN}
The increasing demand for academic document preparation has led to a growing need for efficient template processing systems. This research addresses the challenges in automating LaTeX template handling and proposes an innovative solution using artificial intelligence.

Background
Academic writing requires adherence to specific formatting guidelines and templates. Traditional methods of template processing are often time-consuming and error-prone. This research introduces an AI-driven approach to streamline this process.

Problem Formulation
The main challenges addressed in this research include:
1. Efficient template analysis and processing
2. Accurate content generation following academic standards
3. Multilingual support for template handling

Research Objectives
This study aims to:
1. Develop an automated system for LaTeX template processing
2. Optimize template analysis using machine learning
3. Improve the efficiency of academic document preparation

Research Limitations
The scope of this research is limited to:
1. Academic templates in Indonesian and English
2. LaTeX-based document formats
3. Specific academic document types

Benefits of Research
This research provides:
1. Improved efficiency in academic document preparation
2. Reduced errors in template processing
3. Enhanced accessibility for academic writers
4. Time savings for institutions and researchers

\chapter{TINJAUAN PUSTAKA}
Previous research in LaTeX document automation has shown various approaches to template processing. Smith et al. (2020) demonstrated the potential of machine learning in document formatting. Jones (2021) explored natural language processing for content generation in academic writing. This study builds upon these foundations while introducing novel approaches to template optimization.

% Original content (commented out):
% \documentclass{Base}
\usepackage{framed}
\usepackage[outerbars]{changebar}
\begin{document}





\chapter{The Binomial Model}\label{bin}

In this chapter we will study, in some detail, the simplest
possible nontrivial model of a financial market---the binomial model. This
is a discrete time model, but despite the fact that the main purpose of the
book concerns continuous time models, the binomial model is well worth
studying. The model is very easy to understand, almost all important
concepts which we will study later on already  appear in the binomial case,
the mathematics required to analyze it is at high school level, and last but
not least the binomial model is often used in practice.

%  SECTION   %

\section{The One Period Model} \label{opm}

We start with the one period version of the model. In the next
section we will  (easily) extend the model to an arbitrary number of
periods.

\subsection{Model Description}\label{md}

Running time is denoted by the letter $t$, and by definition we have two
points in time, $t=0$ (``today'') and $t=1$ (``tomorrow''). In the model we
have two assets: a {\bf bond} and a {\bf stock}. At time $t$ the price of a
bond is denoted by $B_t$, and the price of one share of the stock is denoted
by $S_t$. Thus we have two price processes $B$ and~$S$.
\begin{enumerate}[2.]
\item[1.] The bond price process is deterministic and given by
\begin{enumerate}[b.]
\item[a.]This
is a discrete time model, but despite the fact that the main purpose of the
book concerns continuous time models, the binomial model is well worth
studying.
\item[b.]The model is very easy to understand, almost all important
concepts which we will study later on already  appear in the binomial case,
the mathematics required to analyze it is at high school level, and last but
not least the binomial model is often used in practice.
\end{enumerate}
\item[2.]
The constant $R$ is the spot rate for the period, and we can also interpret
the existence of the bond as the existence of a bank with $R$ as its rate of
interest.
\end{enumerate}
The stock price process is a stochastic process, and its dynamical behavior
is described as follows:
\begin{eqnarray}
S_0&=&s,\\
S_1&=&\left\{\begin{array}{@{}c@{}c}
s \cdot u, &\quad \mbox{with probability $p_u$}.\\
s \cdot d, &\quad \mbox{with probability $p_d$}.
\end{array}
\right.
\end{eqnarray}

It is often convenient to write this as
\begin{equation*}
\left\{\begin{array}{@{}r@{}c@{}l}
S_0&\,{=}\,&s,\\
S_1&\,{=}\,&s \cdot Z,
\end{array}
\right.
\end{equation*}
where $Z$ is a stochastic variable defined as\footnote{We assume that today's stock price $s$ is known,
constants $u$, $d$, $p_u$ and $p_d$. We assume that $d<u$, and we have of
course $p_u+p_d=1$.}
\begin{equation*}
Z=\left\{\begin{array}{@{}c@{}c}
u,&\quad \mbox{with probability }p_u.\\
d,&\quad \mbox{with probability }p_d.
\end{array}
\right.
\end{equation*}
\begin{itemize}
\item We\footnote{We assume that $d<u$, and we have of\\[6pt]
\centerline{$a+b=c$}\\[6pt]
course $p_u+p_d=1$.} assume that today's stock price $s$ is known, as are the positive
constants $u$, $d$, $p_u$ and $p_d$. We assume that $d<u$, and we have of
course $p_u+p_d=1$. We can illustrate the price dynamics using the tree structure
in Fig.~\ref{binf1}.
\item We will study the behavior of various {\bf portfolios}
 on the $(B,S)$
market, and to this end we define a portfolio as a vector $h=(x,y)$.
\end{itemize}
%%Figure 1

%\subsubsection{Portfolios and Arbitrage}\label{paa}



We will study\marginpar{\textbf{For Margin}
Everyone wants to make a profit by trading on the market, and in
this context a so called arbitrage portfolio is a dream come true;
this is one of the central concepts of the theory.
} the behavior of various {\bf portfolios}
 on the $(B,S)$
market, and to this end we define a portfolio as a vector $h=(x,y)$. The
interpretation is that $x$ is the number of bonds we hold in our portfolio,
whereas $y$ is the number of units of the stock held by us. Note that it is
quite acceptable for $x$ and $y$ to be positive as well as negative. If, for example,
$x=3$, this means that we have bought three bonds at time $t=0$. If on the
other hand $y=-2$, this means that we have sold two shares of the stock at
time $t=0$. In financial jargon we have a~{\bf long} position in the bond
and a {\bf short} position in the stock. It is an important assumption of
the model that short positions are allowed.

Consider now a fixed portfolio $h=(x,y)$. This portfolio has a deterministic
market value at $t=0$ and a stochastic value at $t=1$.

Everyone wants to make a profit by trading on the market, and in this
context a so called arbitrage portfolio is a dream come true; this is one
of the central concepts of the theory

\begin{definition}
An {\bf arbitrage} portfolio is a portfolio $h$ with the properties
\begin{eqnarray}
V_0^h&=&0,\\
V_1^h&>&0,\quad\mbox{with probability }1.
\end{eqnarray}
\end{definition}

\looseness=-1 An arbitrage portfolio is thus basically a
deterministic money making machine, and we interpret the existence
of an arbitrage portfolio as equivalent to a serious case of
mispricing on the market. It is now natural to investigate when a
given market model is arbitrage free, i.e. when there are no
arbitrage portfolios.

\begin{proposition}
The model above is free of arbitrage if and only if the following conditions
hold:
\begin{equation}\label{noarb}
d \leq (1+R) \leq u.
\end{equation}
\end{proposition}


The condition (\ref{noarb}) has an easy economic interpretation. It simply
says that the return on the stock is not allowed to dominate the return on
the bond and vice versa.

Now assume that (\ref{noarb}) is satisfied. To show that this
implies absence of arbitrage let us consider an arbitrary
portfolio such that $V_0^h=0$. We thus have $x+ys=0$, i.e.
$x=-ys$. Using this relation we can write the value of the
portfolio at $t=1$ as
\begin{equation*}
V_1^h= \left\{\begin{array}{@{}c@{}c}
ys(u-(1+R)),&\quad \mbox{if }Z=u.\\
ys(d-(1+R)), &\quad \mbox{if }Z=d.
\end{array}
\right.
\end{equation*}
Assume now that $y>0$. Then $h$ is an arbitrage strategy if and only if we
have the inequalities
\begin{eqnarray}
u&>&1+R,\\
d&>&1+R,
\end{eqnarray}
but this is impossible because of the condition (\ref{noarb}). The case $y<0$ is treated similarly. \endproof

At first glance this result is perhaps only moderately exciting, but we may
write it in a more suggestive form. To say that (\ref{noarb}) holds is
equivalent to saying that $1+R$ is a convex combination of $u$ and $d$, i.e.
\begin{equation*}
1+R=q_u\cdot u+q_d\cdot d,
\end{equation*}
\looseness=-1 where $q_u,q_d\geq 0$ and $q_u+q_d=1$. In particular we see that the weights
$q_u$ and $q_d$ can be interpreted as probabilities for a new probability
measure $Q$ with the property $Q(Z=u)=q_u$, $Q(Z=d)=q_d$. Denoting
expectation w.r.t. this measure by $E^Q$ we now have the following easy
calculation
\begin{equation*}
\frac 1{1+R} Q {S_1}=11+R(q_usu+q_dsd)=\frac 1{1+R}\cdot
 s(1+R)=s.
\end{equation*}

We thus have the relation
\begin{equation*}
s=\frac 1{1+R} S_1,
\end{equation*}
which to an economist is a well-known relation. It is in fact a {\bf risk
neutral} valuation formula, in the sense that
it gives today's stock price
as the discounted expected value of tomorrow's stock price. Of course
we do not assume that the agents in our market are risk neutral---what we
have shown is only that if we use the $Q$-probabilities instead of the
objective probabilities then we have in fact a risk neutral valuation of the
stock (given absence of arbitrage). A probability measure with this property
is called a {\bf risk neutral measure}, or alternatively a~{\bf risk adjusted measure} or a
{\bf martingale measure}. Martingale measures will play a dominant role in the
sequel so we give a formal definition.

\begin{definition}
A probability measure $Q$ is called a {\bf martingale measure}
if the
following condition holds:
\begin{equation*}
S_0=\frac{1}{1+R} S_1.
\end{equation*}
\end{definition}

We may now state the condition of no arbitrage in the following way.

For the binomial model it is easy to calculate the martingale probabilities.
The proof is left to the reader.

\subsection{Contingent Claims}
\label{bcc}

Let us now assume that the market in the preceding section is arbitrage
free. We go on to study pricing problems for contingent claims.

\begin{definition}
A {\bf contingent claim}
(financial derivative)
is any stochastic variable $X$ of the form $X=\Phi (Z)$, where $Z$ is the stochastic variable driving the
stock price process above. \end{definition}

We interpret a given claim $X$ as a contract which pays $X$ SEK to the
holder of the contract at time $t=1$. See Fig.~\ref{bin_sin1}, where
the value of the claim at each node is given within the corresponding box.
The function $\Phi$ is called the
{\bf contract function}. A typical example would be a European call option
on the stock with strike price $K$. For this option to be interesting we
assume that $sd <K<su$. If $S_1 >K$ then we use the option, pay $K$ to get
the stock and then sell the stock on the market for $su$, thus making a net
profit of $su-K$. If $S_1<K$ then the option is obviously worthless. In this
example we thus have
\begin{equation*}
X=\left\{\begin{array}{@{}l@{}c}
su-K, &\quad\mbox{if $Z=u$},\\
0, &\quad \mbox{if $Z=d$},
\end{array}
\right.
\end{equation*}
and the contract function is given by
\begin{eqnarray}
\Phi(u)&=&su-K,\\
\Phi(d)&=&0.
\end{eqnarray}

Our main problem is now to determine the ``fair'' price, if such an object
exists at all, for a given contingent claim $X$. If we denote the price of $X$ at time $t$ by $X$, then it can be seen that at time $t=1$ the
problem is easy to solve. In order to avoid arbitrage we must (why?) have
\begin{equation*}
X=X,
\end{equation*}
and the hard part of the problem is to determine $X$. To attack
this problem we make a slight detour.

Since we have assumed absence of arbitrage we know that we cannot make
money out of nothing, but it is interesting to study what we {\bf can}
achieve on the market.

\begin{definition}
A given contingent claim $X$ is said to be {\bf reachable}
if there exists a
portfolio $h$ such that
\begin{equation*}
V_1^h=X,
\end{equation*}
with probability $1$. In that case we say that the portfolio $h$ is a
{\bf hedging} portfolio or a
{\bf replicating} portfolio. If all claims can be
replicated we say that the market is {\bf complete}. \end{definition}


If a certain claim $X$ is reachable with replicating portfolio~$h$, then,
from a~financial point of view, there is no difference between holding the
claim and holding the portfolio. No matter what happens on the stock market,
the value of the claim at time $t=1$ will be exactly equal to the value of
the portfolio at $t=1$. Thus the price of the claim should equal the market
value of the portfolio, and we have the following basic pricing principle.

The word ``reasonable'' above can be given a more precise meaning as in the
following proposition. We leave the proof to the reader.

We see that in a complete market we can in fact price all contingent claims,
so it is of great interest to investigate when a given market is complete.
For the binomial model we have the following result.

\begin{table}[b]
\processtable{Table caption }{%
\begin{tabular}{@{}ll@{}}
\toprule
Possible cuts & 1--100\,m\\
\midrule
Total plasma mass & $10-10-2\,{\rm gm}$\\
Ion concentration & $10-10\,{\rm m}\,{-3}$\\
Temperature & 1--40\,keV\\
Pressure & 0.1--5 atmospheres\\
Ion thermal velocity & $100\text{--}1000\,{\rm km}\,{\rm s}$\\
Electron thermal velocity & $0.01c\text{--}0.1c$\\
Magnetic field & 1--10\,T\\
Total plasma current & 0.1--7\,MA\\
\botrule
\end{tabular}}{Table footnote}
\end{table}

\begin{proof}
We fix an arbitrary claim $X$ with contract function $\Phi$, and we want to
show that there exists a portfolio $h=(x,y)$.
\end{proof}

\subsection{Risk Neutral Valuation}
\label{brnv}

Since the binomial model is shown to be complete we can now price any
contingent claim.

\begin{proposition} \label{rnv1}
If the binomial model is free of arbitrage, then the
arbitrage free price of a contingent claim $X$ is given by
\begin{equation} \label{rnv2}
X=\frac{1}{1+R} X.
\end{equation}
Here the martingale measure $Q$ is uniquely determined by the relation
\begin{equation}
\label{rnv3} S_0=\frac{1}{1+R},
\end{equation}
and the explicit expression for $q_u$ and $q_d$ are given in Proposition~\ref{opm100}. Furthermore the claim can be replicated using the portfolio
\begin{eqnarray*}
x&=&\frac{1}{1+R}\cdot \frac{u\Phi (d) -d \Phi (u)}{u-d},\\[5pt]
y&=&\frac{1}{s}\cdot \frac{\Phi (u) - \Phi (d)}{u-d}.
\end{eqnarray*}
\end{proposition}

We see that the formula (\ref{rnv2}) is a ``risk neutral'' valuation
formula, and that the probabilities which are used are just those for which
the stock itself admits a risk neutral valuation. The main economic moral
can now be summarized.

We end by studying a concrete example.

\begin{example}
\emph{We set $s=100$, $u=1.2$, $d=0.8$, $p_u=0.6$, $p_d=0.4$ and, for
computational simplicity, $R=0$. By convention, the monetary unit is the US dollar.
Thus we have the price dynamics}
\begin{eqnarray}
S_0&=&100,\\
S_1&=&
\left\{\begin{array}{@{}r@{}c}
120,&\quad \mbox{{\rm with probability }}0.6.\\
80,&\quad \mbox{{\rm with probability }}0.4.
\end{array}
\right.
\end{eqnarray}
\end{example}

\noindent
If we compute the discounted expected
value (under the objective probability measure $P$) of tomorrow's price  we get
assume a constant deterministic short rate of interest $R$, which is
interpreted as the simple period rate. This means that the bond price
dynamics are given by
\begin{eqnarray}
B_{n+1}&=&(1+R)B_n,\\
B_0&=&1.
\end{eqnarray}

We now go on to define the concept of a dynamic portfolio strategy.

\begin{definition}
A {\bf portfolio strategy} is a stochastic process
such that $h_t$ is a function of $S_0,S_1, \ldots,S_{t-1}$. For a given
portfolio strategy $h$ we set $h_0=h_1$ by convention. The
{\bf value process} corresponding to the portfolio $h$ is defined by
\begin{equation*}
V_t^h=x_t(1+R)+y_tS_t.
\end{equation*}
\end{definition}
The condition above is in fact also sufficient for absence of arbitrage, but
this fact is a little harder to show, and we will prove it later. In any
case we assume that the condition holds.
\begin{lemma}
If the model is free of arbitrage then the following conditions necessarily
must hold.
\begin{equation}
\label{noarb2} d \leq (1+R) \leq u.
\end{equation}
\end{lemma}
It is possible, and not very hard, to give a formal proof of the
proposition, using mathematical induction.
\begin{proposition}
The multiperiod binomial model is complete, i.e. every claim can be
replicated by  a self-financing portfolio.
\end{proposition}
The formal proof will, however,
look rather messy with lots of indices, so instead we prove the proposition
for a concrete example, using a binomial tree.
\begin{example} \label{binexample}
\emph{We set $T=3$, $S_0=80$, $u=1.5$, $d=0.5$, $p_u=0.6$, $p_d=0.4$ and, for
computational simplicity, $R=0$.}
\end{example}

\begin{figure}\includegraphics{fig1}
\caption{Price dynamics to this end we define a portfolio as a vector
$h=(x,y)$} \label{binf1}
\end{figure}


\section*{Unlist}

\begin{unlist}
\item The bond price process is deterministic and given by

\item
The constant $R$ is the spot rate for the period, and we can also
interpret the existence of the bond as the existence of a bank
with $R$ as its rate of interest.
\end{unlist}
The constant $R$ is the spot rate for the period, and we can also
interpret the existence of the bond as the existence of a bank
with $R$ as its rate of interest.

\section*{Description}

\begin{description}
\item The constant $R$ is the spot rate for the period, and we can also
interpret the existence of the bond as the existence of a bank
with $R$ as its rate of interest.


\item
The constant $R$ is the spot rate for the period, and we can also
interpret the existence of the bond as the existence of a bank
with $R$ as its rate of interest.
\end{description}

The constant $R$ is the spot rate for the period, and we can also
interpret the existence of the bond as the existence of a bank
with $R$ as its rate of interest.

\section*{Verse}

\begin{verse}
\item The constant $R$ is the spot rate for the period, and we can also
interpret the existence of the bond as the existence of a bank
with $R$ as its rate of interest.


\item
The constant $R$ is the spot rate for the period, and we can also
interpret the existence of the bond as the existence of a bank
with $R$ as its rate of interest.
\end{verse}


The constant $R$ is the spot rate for the period, and we can also
interpret the existence of the bond as the existence of a bank
with $R$ as its rate of interest.

\section*{Quotation}

\begin{quotation}
\item The constant $R$ is the spot rate for the period, and we can also
interpret the existence of the bond as the existence of a bank
with $R$ as its rate of interest.


\item
The constant $R$ is the spot rate for the period, and we can also
interpret the existence of the bond as the existence of a bank
with $R$ as its rate of interest.
\end{quotation}


The constant $R$ is the spot rate for the period, and we can also
interpret the existence of the bond as the existence of a bank
with $R$ as its rate of interest.

\section*{Quote}

\begin{quote}
\item The constant $R$ is the spot rate for the period, and we can also
interpret the existence of the bond as the existence of a bank
with $R$ as its rate of interest.


\item
The constant $R$ is the spot rate for the period, and we can also
interpret the existence of the bond as the existence of a bank
with $R$ as its rate of interest.
\end{quote}
The constant $R$ is the spot rate for the period, and we can also
interpret the existence of the bond as the existence of a bank
with $R$ as its rate of interest.

\begin{theorem}
The model above is free of arbitrage if and only if the following
conditions hold:
\begin{equation}\label{noarb}
d \leq (1+R) \leq u.
\end{equation}
\end{theorem}

The constant $R$ is the spot rate for the period, and we can also
interpret the existence of the bond as the existence of a bank
with $R$ as its rate of interest.

\begin{fp}
The model above is free of arbitrage if and only if the following
conditions hold:
\begin{equation}\label{noarb}
d \leq (1+R) \leq u.
\end{equation}
\end{fp}

The constant $R$ is the spot rate for the period, and we can also
interpret the existence of the bond as the existence of a bank
with $R$ as its rate of interest.

\begin{corollary}
The model above is free of arbitrage if and only if the following
conditions hold:
\begin{equation}\label{noarb}
d \leq (1+R) \leq u.
\end{equation}
\end{corollary}


The constant $R$ is the spot rate for the period, and we can also
interpret the existence of the bond as the existence of a bank
with $R$ as its rate of interest.

\begin{assumption}
The model above is free of arbitrage if and only if the following
conditions hold:
\begin{equation}\label{noarb}
d \leq (1+R) \leq u.
\end{equation}
\end{assumption}

The constant $R$ is the spot rate for the period, and we can also
interpret the existence of the bond as the existence of a bank
with $R$ as its rate of interest.

\begin{condition}
The model above is free of arbitrage if and only if the following
conditions hold:
\begin{equation}\label{noarb}
d \leq (1+R) \leq u.
\end{equation}
\end{condition}


The constant $R$ is the spot rate for the period, and we can also
interpret the existence of the bond as the existence of a bank
with $R$ as its rate of interest.

\begin{idea}
The model above is free of arbitrage if and only if the following
conditions hold:
\begin{equation}\label{noarb}
d \leq (1+R) \leq u.
\end{equation}
\end{idea}

The constant $R$ is the spot rate for the period, and we can also
interpret the existence of the bond as the existence of a bank
with $R$ as its rate of interest.

\begin{remark}
The model above is free of arbitrage if and only if the following
conditions hold:
\begin{equation}\label{noarb}
d \leq (1+R) \leq u.
\end{equation}
\end{remark}

The constant $R$ is the spot rate for the period, and we can also
interpret the existence of the bond as the existence of a bank
with $R$ as its rate of interest.

\begin{result}
The model above is free of arbitrage if and only if the following
conditions hold:
\begin{equation}\label{noarb}
d \leq (1+R) \leq u.
\end{equation}
\end{result}
The constant $R$ is the spot rate for the period, and we can also
interpret the existence of the bond as the existence of a bank
with $R$ as its rate of interest.

\begin{nota}
The model above is free of arbitrage if and only if the following
conditions hold:
\begin{equation}\label{noarb}
d \leq (1+R) \leq u.
\end{equation}
\end{nota}

\begin{Exercises}
The constant $R$ is the spot rate for the period, and we can
also interpret the existence of the bond as the existence of a
bank with $R$ as its rate of interest.
\end{Exercises}

\section*{Boxed Text Crossed the Page}

\begin{mibox}

\brule

\cbstart

\vspace*{3pt}

\boxhead{{INVESTIGATION}}

\vspace*{-5.5pt}
%\invconttrue\begin{investigation}{
\boxsubhead{Blind Embedding and Linear Correlation Detection}

The embedding algorithm in the system we describe here implements
a blind embedder. We denote this algorithm by EBLIND, which refers
to this {\textit{specific}} example of blind embedding rather than
the generic concept of blind embedding. In fact, there are many
other algorithms for blind embedding.

\boxindent The detection algorithm uses linear correlation as its
detection metric. This is a very common detection metric, which is
discussed further in Section~3.5.

\boxindent To keep things simple, we code only one bit of
information. Thus, {m} is either\break 1 or 0. We assume that we
are working with only grayscale images. Most of the algorithms
presented in this book share these simplifications. Methods of
encoding more than one bit are discussed in Chapters~4 and~5.

\vspace*{12pt} \setcounter{figure}{9}

%\begin{figure}%9
\centerline{\fbox{\hbox to 6pc{\vbox to 4pc{\vfill}}}}
\boxfigcap{Distribution of linear correlations between images and
a low-pass filtered random noise pattern.\label{fig:dlc-pink}}
%\end{figure}
\vspace*{8pt}

watermarked image in Figure~3.8 has significantly worse fidelity
than that in Figure~3.7, because the human eye is more sensitive
to low-frequency patterns than to high-frequency patterns.

If a certain claim $X$ is reachable with replicating
portfolio~$h$, then, from a~financial point of view, there is no
difference between holding the claim and holding the portfolio. No
matter what happens on the stock market, the value of the claim at
time $t=1$ will be exactly equal to the value of the portfolio at
$t=1$. Thus the price of the claim should equal the market value
of the portfolio, and we have the following basic pricing
principle.

If a certain claim $X$ is reachable with replicating
portfolio~$h$, then, from a~financial point of view, there is no
difference between holding the claim and holding the portfolio. No
matter what happens on the stock market, the value of the claim at
time $t=1$ will be exactly equal to the value of the portfolio at
$t=1$. Thus the price of the claim should equal the market value
of the portfolio, and we have the following basic pricing
principle.

If a certain claim $X$ is reachable with replicating
portfolio~$h$, then, from a~financial point of view, there is no
difference between holding the claim and holding the portfolio.

\vspace*{-1.5pt}

\brule

\cbend


\end{mibox}

No
matter what happens on the stock market, the value of the claim at
time $t=1$ will be exactly equal to the value of the portfolio at
$t=1$. Thus the price of the claim should equal the market value
of the portfolio, and we have the following basic pricing
principle.

The reader may wish to view (or revisit) the animation
``Inductive and Radiative Coupling'' on the CD.

\boxindent Mathematica is very flexible and most calculations can be carried
out in more than one way.  Depending on how you think, some
sequences of calculations may make more sense to you than others,
even if they are less efficient than the most efficient way to
perform the desired calculations.

Often, the difference in time
required for Mathematica to perform equivalent\,--\,but
different\,--\,calculations is quite small.  For the beginner, we
think it is wisest to work with familiar calculations first and
then efficiency.

\begin{exampled}
Calculate (a)  $121+542$; (b) $3231-9876$; (c)  $(-23)(76)$; (d)
$(22341)(832748)$ $(387281)$; and (e) $\dfrac{467}{31}$.
\end{exampled}

\begin{solution}
These calculations are carried out in the following screen shot.
In each case, the input is typed  and then evaluated by pressing
\textbf{Enter}.  In the last case, the \textbf{Basic Math}
template is used to enter the fraction.
\end{solution}

Mathematica is very flexible and most calculations can be carried
out in more than one way.  Depending on how you think, some
sequences of calculations may make more sense to you than others,
even if they are less efficient than the most efficient way to
perform the desired calculations.  Often, the difference in time
required for Mathematica to perform equivalent\,--\,but
different\,--\,calculations is quite small.  For the beginner, we
think it is wisest to work with familiar calculations first and
then efficiency.

\begin{thenumbibliography}{99}
\bibitem{AJ}
Amin, K. \& Jarrow, R.  (1991).  Pricing Foreign Currency Options under
Stochastic Interest Rates.
{\em Journal of International Money and Finance} {\bf 10}, 310--329.
%
\bibitem{ABD}
Anderson, N., Breedon, F., Deacon, M., Derry, A. \& Murphy, G. (1996).
{\em Estimating and Interpreting the Yield Curve.} Wiley, Chichester.
%
\bibitem{AD}
Artzner, P. \& Delbaen, F.  (1989). Term Structure of Interest Rates: The Martingale Approach.
{\em Advances in Applied Mathematics} {\bf 10}, 95--129.

\bibitem{AJ}
Amin, K. \& Jarrow, R.  (1991).  Pricing Foreign Currency Options
under Stochastic Interest Rates. {\em Journal of International
Money and Finance} {\bf 10}, 310--329.
%
\bibitem{ABD}
Anderson, N., Breedon, F., Deacon, M., Derry, A. \& Murphy, G.
(1996). {\em Estimating and Interpreting the Yield Curve.} Wiley,
Chichester.
%
\bibitem{AD}
Artzner, P. \& Delbaen, F.  (1989). Term Structure of Interest
Rates: The Martingale Approach. {\em Advances in Applied
Mathematics} {\bf 10}, 95--129.

\bibitem{AJ}
Amin, K. \& Jarrow, R.  (1991).  Pricing Foreign Currency Options
under Stochastic Interest Rates. {\em Journal of International
Money and Finance} {\bf 10}, 310--329.
%
\bibitem{ABD}
Anderson, N., Breedon, F., Deacon, M., Derry, A. \& Murphy, G.
(1996). {\em Estimating and Interpreting the Yield Curve.} Wiley,
Chichester.
%
\bibitem{AD}
Artzner, P. \& Delbaen, F.  (1989). Term Structure of Interest
Rates: The Martingale Approach. {\em Advances in Applied
Mathematics} {\bf 10}, 95--129.

\bibitem{AJ}
Amin, K. \& Jarrow, R.  (1991).  Pricing Foreign Currency Options
under Stochastic Interest Rates. {\em Journal of International
Money and Finance} {\bf 10}, 310--329.
%
\bibitem{ABD}
Anderson, N., Breedon, F., Deacon, M., Derry, A. \& Murphy, G.
(1996). {\em Estimating and Interpreting the Yield Curve.} Wiley,
Chichester.
%
\bibitem{AD}
Artzner, P. \& Delbaen, F.  (1989). Term Structure of Interest
Rates: The Martingale Approach. {\em Advances in Applied
Mathematics} {\bf 10}, 95--129.

\bibitem{AJ}
Amin, K. \& Jarrow, R.  (1991).  Pricing Foreign Currency Options
under Stochastic Interest Rates. {\em Journal of International
Money and Finance} {\bf 10}, 310--329.
%
\bibitem{ABD}
Anderson, N., Breedon, F., Deacon, M., Derry, A. \& Murphy, G.
(1996). {\em Estimating and Interpreting the Yield Curve.} Wiley,
Chichester.
%
\bibitem{AD}
Artzner, P. \& Delbaen, F.  (1989). Term Structure of Interest
Rates: The Martingale Approach. {\em Advances in Applied
Mathematics} {\bf 10}, 95--129.

\bibitem{AJ}
Amin, K. \& Jarrow, R.  (1991).  Pricing Foreign Currency Options
under Stochastic Interest Rates. {\em Journal of International
Money and Finance} {\bf 10}, 310--329.
%
\bibitem{ABD}
Anderson, N., Breedon, F., Deacon, M., Derry, A. \& Murphy, G.
(1996). {\em Estimating and Interpreting the Yield Curve.} Wiley,
Chichester.
%
\bibitem{AD}
Artzner, P. \& Delbaen, F.  (1989). Term Structure of Interest
Rates: The Martingale Approach. {\em Advances in Applied
Mathematics} {\bf 10}, 95--129.


\bibitem{AJ}
Amin, K. \& Jarrow, R.  (1991).  Pricing Foreign Currency Options under
Stochastic Interest Rates.
{\em Journal of International Money and Finance} {\bf 10}, 310--329.
%
\bibitem{ABD}
Anderson, N., Breedon, F., Deacon, M., Derry, A. \& Murphy, G. (1996).
{\em Estimating and Interpreting the Yield Curve.} Wiley, Chichester.
%
\bibitem{AD}
Artzner, P. \& Delbaen, F.  (1989). Term Structure of Interest Rates: The Martingale Approach.
{\em Advances in Applied Mathematics} {\bf 10}, 95--129.

\bibitem{AJ}
Amin, K. \& Jarrow, R.  (1991).  Pricing Foreign Currency Options
under Stochastic Interest Rates. {\em Journal of International
Money and Finance} {\bf 10}, 310--329.
%
\bibitem{ABD}
Anderson, N., Breedon, F., Deacon, M., Derry, A. \& Murphy, G.
(1996). {\em Estimating and Interpreting the Yield Curve.} Wiley,
Chichester.
%
\bibitem{AD}
Artzner, P. \& Delbaen, F.  (1989). Term Structure of Interest
Rates: The Martingale Approach. {\em Advances in Applied
Mathematics} {\bf 10}, 95--129.

\end{thenumbibliography}

\begin{theunbibliography}{99}
\bibitem{AJ}
Amin, K. \& Jarrow, R.  (1991).  Pricing Foreign Currency Options under
Stochastic Interest Rates.
{\em Journal of International Money and Finance} {\bf 10}, 310--329.
%
\bibitem{ABD}
Anderson, N., Breedon, F., Deacon, M., Derry, A. \& Murphy, G. (1996).
{\em Estimating and Interpreting the Yield Curve.} Wiley, Chichester.
%
\bibitem{AD}
Artzner, P. \& Delbaen, F.  (1989). Term Structure of Interest Rates: The Martingale Approach.
{\em Advances in Applied Mathematics} {\bf 10}, 95--129.
\end{theunbibliography}


\clearpage

\setcounter{chapter}{0}

\appendix

\chapter{}

We will study the behavior of various {\bf portfolios}
 on the $(B,S)$
market, and to this end we define a portfolio as a vector
$h=(x,y)$. The interpretation is that $x$ is the number of bonds
we hold in our portfolio, whereas $y$ is the number of units of
the stock held by us. Note that it is quite acceptable for $x$ and
$y$ to be positive as well as negative. If, for example, $x=3$,
this means that we have bought three bonds at time $t=0$. If on
the other hand $y=-2$, this means that we have sold two shares of
the stock at time $t=0$. In financial jargon we have a~{\bf long}
position in the bond and a {\bf short} position in the stock. It
is an important assumption of the model that short positions are
allowed.

\section{Appendix Head1}

Consider now a fixed portfolio $h=(x,y)$. This portfolio has a
deterministic market value at $t=0$ and a stochastic value at
$t=1$.

Everyone wants to make a profit by trading on the market, and in
this context a so called arbitrage portfolio is a dream come true;
this is one of the central concepts of the theory.

We will study the behavior of various {\bf portfolios}
 on the $(B,S)$
market, and to this end we define a portfolio as a vector
$h=(x,y)$. The interpretation is that $x$ is the number of bonds
we hold in our portfolio, whereas $y$ is the number of units of
the stock held by us. Note that it is quite acceptable for $x$ and
$y$ to be positive as well as negative. If, for example, $x=3$,
this means that we have bought three bonds at time $t=0$. If on
the other hand $y=-2$, this means that we have sold two shares of
the stock at time $t=0$. In financial jargon we have a~{\bf long}
position in the bond and a {\bf short} position in the stock. It
is an important assumption of the model that short positions are
allowed.

\section{Appendix Head2}

Consider now a fixed portfolio $h=(x,y)$. This portfolio has a
deterministic market value at $t=0$ and a stochastic value at
$t=1$.

Everyone wants to make a profit by trading on the market, and in
this context a so called arbitrage portfolio is a dream come true;
this is one of the central concepts of the theory.

We will study the behavior of various {\bf portfolios}
 on the $(B,S)$
market, and to this end we define a portfolio as a vector
$h=(x,y)$. The interpretation is that $x$ is the number of bonds
we hold in our portfolio, whereas $y$ is the number of units of
the stock held by us. Note that it is quite acceptable for $x$ and
$y$ to be positive as well as negative. If, for example, $x=3$,
this means that we have bought three bonds at time $t=0$. If on
the other hand $y=-2$, this means that we have sold two shares of
the stock at time $t=0$. In financial jargon we have a~{\bf long}
position in the bond and a {\bf short} position in the stock. It
is an important assumption of the model that short positions are
allowed.

\chapter{}

Consider now a fixed portfolio $h=(x,y)$. This portfolio has a
deterministic market value at $t=0$ and a stochastic value at
$t=1$.

\section{Appendix Head1}

Everyone wants to make a profit by trading on the market, and in
this context a so called arbitrage portfolio is a dream come true;
this is one of the central concepts of the theory.

We will study the behavior of various {\bf portfolios}
 on the $(B,S)$
market, and to this end we define a portfolio as a vector
$h=(x,y)$. The interpretation is that $x$ is the number of bonds
we hold in our portfolio, whereas $y$ is the number of units of
the stock held by us. Note that it is quite acceptable for $x$ and
$y$ to be positive as well as negative. If, for example, $x=3$,
this means that we have bought three bonds at time $t=0$. If on
the other hand $y=-2$, this means that we have sold two shares of
the stock at time $t=0$. In financial jargon we have a~{\bf long}
position in the bond and a {\bf short} position in the stock. It
is an important assumption of the model that short positions are
allowed.

\section{Appendix Head2}

Consider now a fixed portfolio $h=(x,y)$. This portfolio has a
deterministic market value at $t=0$ and a stochastic value at
$t=1$.

Everyone wants to make a profit by trading on the market, and in
this context a so called arbitrage portfolio is a dream come true;
this is one of the central concepts of the theory.

We will study the behavior of various {\bf portfolios}
 on the $(B,S)$
market, and to this end we define a portfolio as a vector
$h=(x,y)$. The interpretation is that $x$ is the number of bonds
we hold in our portfolio, whereas $y$ is the number of units of
the stock held by us. Note that it is quite acceptable for $x$ and
$y$ to be positive as well as negative. If, for example, $x=3$,
this means that we have bought three bonds at time $t=0$. If on
the other hand $y=-2$, this means that we have sold two shares of
the stock at time $t=0$. In financial jargon we have a~{\bf long}
position in the bond and a {\bf short} position in the stock. It
is an important assumption of the model that short positions are
allowed.

Consider now a fixed portfolio $h=(x,y)$. This portfolio has a
deterministic market value at $t=0$ and a stochastic value at
$t=1$.

Everyone wants to make a profit by trading on the market, and in
this context a so called arbitrage portfolio is a dream come true;
this is one of the central concepts of the theory.

We will study the behavior of various {\bf portfolios}
 on the $(B,S)$
market, and to this end we define a portfolio as a vector
$h=(x,y)$. The interpretation is that $x$ is the number of bonds
we hold in our portfolio, whereas $y$ is the number of units of
the stock held by us. Note that it is quite acceptable for $x$ and
$y$ to be positive as well as negative. If, for example, $x=3$,
this means that we have bought three bonds at time $t=0$. If on
the other hand $y=-2$, this means that we have sold two shares of
the stock at time $t=0$. In financial jargon we have a~{\bf long}
position in the bond and a {\bf short} position in the stock. It
is an important assumption of the model that short positions are
allowed.

Consider now a fixed portfolio $h=(x,y)$. This portfolio has a
deterministic market value at $t=0$ and a stochastic value at
$t=1$.

Everyone wants to make a profit by trading on the market, and in
this context a so called arbitrage portfolio is a dream come true;
this is one of the central concepts of the theory.


\end{document}


\end{document}